\documentclass[10pt,a4paper]{amsart}
\usepackage{amsmath,amssymb,amsthm,graphicx}
% Retained minimal preamble for the Stacks Project chapter more-algebra.tex (More on Algebra).
% The chapter file pinned-original-source.tex carries only \input{preamble} on its first line, so no
% preamble text lives inside the chapter itself.  The Stacks Project uses one single shared preamble for
% every chapter; the snapshot stages it identically as sources/stacks-chow/preamble.tex and
% sources/stacks-etale-cohomology/preamble.tex.  The block below is copied mechanically, line for line,
% from sources/stacks-chow/preamble.tex lines 14-18 (the Xy-pic setup) and 154-217 (the native theorem
% declarations together with the mathematical macros).  Nothing was added, removed or redefined; the
% wrapper class and the amsmath/amssymb/amsthm/graphicx base packages are supplied by the assembler.
% Chapter: more-algebra.tex -- More on Algebra.
% For commutative diagrams we use Xy-pic
\usepackage[all]{xy}

% We use 2cell for 2-commutative diagrams.
\xyoption{2cell}

% Theorem environments.
%
\theoremstyle{plain}
\newtheorem{theorem}[subsection]{Theorem}
\newtheorem{proposition}[subsection]{Proposition}
\newtheorem{lemma}[subsection]{Lemma}

\theoremstyle{definition}
\newtheorem{definition}[subsection]{Definition}
\newtheorem{example}[subsection]{Example}
\newtheorem{exercise}[subsection]{Exercise}
\newtheorem{situation}[subsection]{Situation}

\theoremstyle{remark}
\newtheorem{remark}[subsection]{Remark}
\newtheorem{remarks}[subsection]{Remarks}

\numberwithin{equation}{subsection}

% Macros
%
\def\lim{\mathop{\mathrm{lim}}\nolimits}
\def\colim{\mathop{\mathrm{colim}}\nolimits}
\def\Spec{\mathop{\mathrm{Spec}}}
\def\Hom{\mathop{\mathrm{Hom}}\nolimits}
\def\Ext{\mathop{\mathrm{Ext}}\nolimits}
\def\SheafHom{\mathop{\mathcal{H}\!\mathit{om}}\nolimits}
\def\SheafExt{\mathop{\mathcal{E}\!\mathit{xt}}\nolimits}
\def\Sch{\mathit{Sch}}
\def\Mor{\mathop{\mathrm{Mor}}\nolimits}
\def\Ob{\mathop{\mathrm{Ob}}\nolimits}
\def\Sh{\mathop{\mathit{Sh}}\nolimits}
\def\NL{\mathop{N\!L}\nolimits}
\def\CH{\mathop{\mathrm{CH}}\nolimits}
\def\proetale{{pro\text{-}\acute{e}tale}}
\def\etale{{\acute{e}tale}}
\def\QCoh{\mathit{QCoh}}
\def\Ker{\mathop{\mathrm{Ker}}}
\def\Im{\mathop{\mathrm{Im}}}
\def\Coker{\mathop{\mathrm{Coker}}}
\def\Coim{\mathop{\mathrm{Coim}}}

% Boxtimes
%
\DeclareMathSymbol{\boxtimes}{\mathbin}{AMSa}{"02}

%
% Macros for moduli stacks/spaces
%
\def\QCohstack{\mathcal{QC}\!\mathit{oh}}
\def\Cohstack{\mathcal{C}\!\mathit{oh}}
\def\Spacesstack{\mathcal{S}\!\mathit{paces}}
\def\Quotfunctor{\mathrm{Quot}}
\def\Hilbfunctor{\mathrm{Hilb}}
\def\Curvesstack{\mathcal{C}\!\mathit{urves}}
\def\Polarizedstack{\mathcal{P}\!\mathit{olarized}}
\def\Complexesstack{\mathcal{C}\!\mathit{omplexes}}
% \Pic is the operator that assigns to X its picard group, usage \Pic(X)
% \Picardstack_{X/B} denotes the Picard stack of X over B
% \Picardfunctor_{X/B} denotes the Picard functor of X over B
\def\Pic{\mathop{\mathrm{Pic}}\nolimits}
\def\Picardstack{\mathcal{P}\!\mathit{ic}}
\def\Picardfunctor{\mathrm{Pic}}
\def\Deformationcategory{\mathcal{D}\!\mathit{ef}}

% --- appended 2026-10-12 (Chief): the project own definitions of the reference/slogan/history
% display environments, copied mechanically from the shared Stacks preamble lines 8-12
% (\usepackage{verbatim} plus the three \newenvironment lines, which the chapter preamble keeps
% verbatim from the project preamble). Nothing added, removed or redefined: these environments
% carry no mathematical content and the project itself typesets them as comments. Without these
% three lines a counted statement whose own text uses them fails with Environment slogan undefined.
\usepackage{verbatim}
\newenvironment{reference}{\comment}{\endcomment}
\newenvironment{slogan}{\comment}{\endcomment}
\newenvironment{history}{\comment}{\endcomment}
\usepackage{cleveref}
\crefname{originallemma}{original lemma}{original lemmas}
\Crefname{originallemma}{Original lemma}{Original lemmas}
\makeatletter
\newlabel{more-algebra-lemma-lift-idempotent-upstairs}{{07M4}{source}{}{sourcepointer-1}{}}
\newlabel{more-algebra-lemma-lift-idempotent-upstairs@cref}{{[originallemma][0][]07M4}{[1][0][]source}{}{}{}}
\makeatother
\crefname{originallemma}{original lemma}{original lemmas}
\Crefname{originallemma}{Original lemma}{Original lemmas}
\makeatletter
\newlabel{lemma-helper-integral}{{09XG}{source}{}{sourcepointer-2}{}}
\newlabel{lemma-helper-integral@cref}{{[originallemma][0][]09XG}{[1][0][]source}{}{}{}}
\makeatother
\crefname{originallemma}{original lemma}{original lemmas}
\Crefname{originallemma}{Original lemma}{Original lemmas}
\makeatletter
\newlabel{lemma-lift-factorization-easy}{{07M1}{source}{}{sourcepointer-3}{}}
\newlabel{lemma-lift-factorization-easy@cref}{{[originallemma][0][]07M1}{[1][0][]source}{}{}{}}
\makeatother
\crefname{originallemma}{original lemma}{original lemmas}
\Crefname{originallemma}{Original lemma}{Original lemmas}
\makeatletter
\newlabel{algebra-lemma-integral-going-up}{{00GU}{source}{}{sourcepointer-4}{}}
\newlabel{algebra-lemma-integral-going-up@cref}{{[originallemma][0][]00GU}{[1][0][]source}{}{}{}}
\makeatother
\crefname{originallemma}{original lemma}{original lemmas}
\Crefname{originallemma}{Original lemma}{Original lemmas}
\makeatletter
\newlabel{algebra-lemma-going-up-closed}{{00HZ}{source}{}{sourcepointer-5}{}}
\newlabel{algebra-lemma-going-up-closed@cref}{{[originallemma][0][]00HZ}{[1][0][]source}{}{}{}}
\makeatother
\crefname{originallemma}{original lemma}{original lemmas}
\Crefname{originallemma}{Original lemma}{Original lemmas}
\makeatletter
\newlabel{lemma-separate-image-closed-from-closed}{{07M3}{source}{}{sourcepointer-6}{}}
\newlabel{lemma-separate-image-closed-from-closed@cref}{{[originallemma][0][]07M3}{[1][0][]source}{}{}{}}
\makeatother
\crefname{originallemma}{original lemma}{original lemmas}
\Crefname{originallemma}{Original lemma}{Original lemmas}
\makeatletter
\newlabel{algebra-lemma-disjoint-decomposition}{{00EE}{source}{}{sourcepointer-7}{}}
\newlabel{algebra-lemma-disjoint-decomposition@cref}{{[originallemma][0][]00EE}{[1][0][]source}{}{}{}}
\makeatother
\title{Lifting an idempotent modulo an ideal after an etale base change that is an isomorphism on the quotient}
\author{The Stacks Project Authors}
\date{stacks-project:more-algebra.tex:local-snapshot-20261012; GNU FDL 1.2 or later}
\begin{document}
\maketitle
\setcounter{subsection}{1}
\noindent\textit{Editorial note.} All delivered mathematical text is the author's own, copied from the pinned Stacks Project chapter file more-algebra.tex; no statement or proof step was rewritten, reordered or removed.\par\medskip
\noindent\textit{Editorial note.} The chapter file pinned-original-source.tex carries only \textbackslash{}input\{preamble\} on its first line, so no preamble text lives inside the chapter. The retained file sources/stacks-more-algebra/retained-preamble.tex is a mechanical, line-for-line copy of the project's single shared preamble as staged in the snapshot (Xy-pic setup plus theorem declarations and mathematical macros). Nothing was added, removed or redefined, and the excerpt serves only as preamble, never as a span coordinate.\par\medskip
\noindent\textit{Editorial note.} Both counted spans start at the \textbackslash{}begin line of their environment, so the statement and the author's proof are delivered whole; surrounding prose of the section is deliberately outside the counted spans and is carried as standard prerequisites instead.\par\medskip
\noindent\textit{Editorial note.} The citation fields (source\_id, source\_version, authors, source\_url, license.name, source\_title) are identical to those of every other card this lane has produced from this chapter, so the batch selector treats them as one source.\par\medskip
\noindent\textit{Allowed standard prerequisites (editorial).}\begin{itemize}
\item Existence of a monic polynomial modulo the ideal vanishing at a lift of the idempotent. Lemma \textbackslash{}ref\{lemma-helper-integral\}, cited at proof lines 1988-1989. Produces the polynomial f whose factorisation drives the construction.
\item Factorisation of that polynomial after an etale base change inducing an isomorphism modulo the ideal. Lemma \textbackslash{}ref\{lemma-lift-factorization-easy\}, cited at proof lines 1990-1992. Gives the two factors g and h congruent to a variable power and to a power of the variable minus one modulo the ideal.
\item An integral ring map is closed on spectra, and closed disjoint loci can be separated by an element. Algebra, Lemmas \textbackslash{}ref\{algebra-lemma-integral-going-up\} and \textbackslash{}ref\{algebra-lemma-going-up-closed\}, cited at proof lines 2007-2009, and Lemma \textbackslash{}ref\{lemma-separate-image-closed-from-closed\}, cited at proof lines 2010-2012. Produce the element whose localisation makes the two loci cover the spectrum.
\item Idempotents correspond to open and closed subsets, and the associated idempotent is unique. Algebra, Lemma \textbackslash{}ref\{algebra-lemma-disjoint-decomposition\}, cited at proof lines 2014 and 2017. Give the idempotent of the disjoint decomposition and identify it as the required lift.
\end{itemize}\medskip
\ifdefined\hypertarget\hypertarget{sourcepointer-1}{}\fi
\noindent\textit{Source reference.} Original lemma 07M4: focused-seeds/AGP/sources/stacks-snapshot-20261011/tags.txt, source lines 9714--9714. The lemma's own canonical tag, quoted so that readers can resolve the label inside the delivered text; it points at the counted statement itself and carries no proof content.\par\medskip
\ifdefined\hypertarget\hypertarget{sourcepointer-2}{}\fi
\noindent\textit{Source reference.} Original lemma 09XG: focused-seeds/AGP/sources/stacks-snapshot-20261011/tags.txt, source lines 12529--12529. Navigation pointer to a statement the counted proof refers to but does not contain; nothing load-bearing is delegated to it.\par\medskip
\ifdefined\hypertarget\hypertarget{sourcepointer-3}{}\fi
\noindent\textit{Source reference.} Original lemma 07M1: focused-seeds/AGP/sources/stacks-snapshot-20261011/tags.txt, source lines 9711--9711. Navigation pointer to a statement the counted proof refers to but does not contain; nothing load-bearing is delegated to it.\par\medskip
\ifdefined\hypertarget\hypertarget{sourcepointer-4}{}\fi
\noindent\textit{Source reference.} Original lemma 00GU: focused-seeds/AGP/sources/stacks-snapshot-20261011/tags.txt, source lines 620--620. Navigation pointer to a statement the counted proof refers to but does not contain; nothing load-bearing is delegated to it.\par\medskip
\ifdefined\hypertarget\hypertarget{sourcepointer-5}{}\fi
\noindent\textit{Source reference.} Original lemma 00HZ: focused-seeds/AGP/sources/stacks-snapshot-20261011/tags.txt, source lines 661--661. Navigation pointer to a statement the counted proof refers to but does not contain; nothing load-bearing is delegated to it.\par\medskip
\ifdefined\hypertarget\hypertarget{sourcepointer-6}{}\fi
\noindent\textit{Source reference.} Original lemma 07M3: focused-seeds/AGP/sources/stacks-snapshot-20261011/tags.txt, source lines 9713--9713. Navigation pointer to a statement the counted proof refers to but does not contain; nothing load-bearing is delegated to it.\par\medskip
\ifdefined\hypertarget\hypertarget{sourcepointer-7}{}\fi
\noindent\textit{Source reference.} Original lemma 00EE: focused-seeds/AGP/sources/stacks-snapshot-20261011/tags.txt, source lines 531--531. Navigation pointer to a statement the counted proof refers to but does not contain; nothing load-bearing is delegated to it.\par\medskip
\begin{lemma}
\label{lemma-lift-idempotent-upstairs}
Let $A$ be a ring, let $I \subset A$ be an ideal.
Let $A \to B$ be an integral ring map.
Let $\overline{e} \in B/IB$ be an idempotent.
Then there exists an \'etale ring map $A \to A'$
which induces an isomorphism $A/I \to A'/IA'$ and an idempotent
$e' \in B \otimes_A A'$ lifting $\overline{e}$.
\end{lemma}

\begin{proof}
Choose an element $y \in B$ lifting $\overline{e}$.
Choose $f \in A[x]$ as in Lemma \ref{lemma-helper-integral} for $y$.
By Lemma \ref{lemma-lift-factorization-easy}
we can find an \'etale ring map $A \to A'$ which induces
an isomorphism $A/I \to A'/IA'$ and such that $f = gh$
in $A[x]$ with $g(x) = x^d \bmod IA'$ and $h(x) = (x - 1)^d \bmod IA'$.
After replacing $A$ by $A'$ we may assume that the factorization
is defined over $A$. In that case we see that
$b_1 = g(y) \in B$ is a lift of $\overline{e}^d = \overline{e}$ and
$b_2 = h(y) \in B$ is a lift of
$(\overline{e} - 1)^d = (-1)^d (1 - \overline{e})^d = (-1)^d(1 - \overline{e})$
and moreover $b_1b_2 = 0$. Thus $(b_1, b_2)B/IB = B/IB$ and
$V(b_1, b_2) \subset \Spec(B)$ is disjoint from $V(IB)$. Since
$\Spec(B) \to \Spec(A)$ is closed (see
Algebra, Lemmas \ref{algebra-lemma-integral-going-up} and
\ref{algebra-lemma-going-up-closed})
we can find an $a \in A$ which maps to an invertible element
of $A/I$ whose image in $B$ lies in $(b_1, b_2)$, see
Lemma \ref{lemma-separate-image-closed-from-closed}.
After replacing $A$ by the localization $A_a$ we get that
$(b_1, b_2) = B$. Then $\Spec(B) = D(b_1) \amalg D(b_2)$;
disjoint union because $b_1b_2 = 0$ and covers
$\Spec(B)$ because $(b_1, b_2) = B$. Let $e \in B$ be the idempotent
corresponding to the open and closed subset $D(b_1)$, see
Algebra, Lemma \ref{algebra-lemma-disjoint-decomposition}.
Since $b_1$ is a lift of $\overline{e}$ and $b_2$ is a
lift of $\pm (1 - \overline{e})$ we conclude that $e$ is
a lift of $\overline{e}$ by the uniqueness statement in
Algebra, Lemma \ref{algebra-lemma-disjoint-decomposition}.
\end{proof}

\end{document}
